% Level 1: tables, and the Lean they need
\begin{frame}
  \frametitle{A Relation Is a Finite Set of Rows}
  {\centering\small
    \smalltable{$\rel{Member}$}{ll}{$\att{name}$ & $\att{dept}$}
      {Ada & Maths \\ Alan & CS \\ Grace & CS}\hspace{12mm}
    \smalltable{$\rel{Attended}$}{lr}{$\att{name}$ & $\att{week}$}
      {Ada & 1 \\ Alan & 1 \\ Alan & 2 \\ Grace & 3}\hspace{12mm}
    \smalltable{$\rel{Session}$}{rl}{$\att{week}$ & $\att{topic}$}
      {1 & Lists \\ 2 & Proofs \\ 3 & Types}\par}
  \begin{itemize}
    \item A \textbf{schema} is a set of attributes, such as
      $\{\att{name}, \att{dept}\}$; a \textbf{row} gives each a value.
    \item A \textbf{relation} is a finite \emph{set} of rows over one
      schema: no order, no duplicates.
    \item A \textbf{database} gives a relation to each table name.
  \end{itemize}
  \slideref{E. F. Codd. \textbf{A Relational Model of Data for Large Shared Data Banks}. CACM 13(6), 1970.}
  \leveltag{1}
\end{frame}

\leanframe{Level 1 Needs Structures, Inductive Types and Lists}{1}{%
  define a \code{structure} and an \code{inductive} type, and
    pattern-match on them & \rds{fpil-struct, fpil-data} \\
  write recursive functions on lists; use \code{Option} &
    \rds{fpil-poly} \\
  derive \code{BEq}, \code{Repr}, \code{DecidableEq} for your types &
    \rds{fpil-deriving, tpil-dec} \\
  write our example database, and \code{\#eval} it & \rds{fpil-eval} \\
}{Then FPIL 7.3.1--7.3.2 builds column types, schemas, rows and tables.}
